\subsection{基本权、支配权}

假设 R 是约化的。设 C 是 R 的一个室，B 是相应的基。由于 R 是约化的，\(B^{-} = \{\alpha^{-}\}_{\alpha \in B}\) 是 \(R^{-}\) 的一个基。于是 \((\overline{\omega}_{\alpha})_{\alpha \in B}\) 是 \(B^{-}\) 的对偶基，因而是权群的一个基；其元素称为基本权（相对于 B，或相对于 C）；若 B 的元素记为 \((\alpha_{1}, \ldots, \alpha_{l})\)，则相应的基本权记为 \((\overline{\omega}_{1}, \ldots, \overline{\omega}_{l})\)。

设 \(x \in V\)。则 \(x \in C\) 当且仅当 \(\langle x, \alpha \rangle > 0\) 对所有 \(\alpha \in B\) 成立。因此，\(C\) 是由 \(\bar{\omega}_{\alpha}\) 作线性组合所成的集合，且系数为 \(\geqslant 0\)。

Cartan 矩阵的元素 \(n(\alpha, \beta) = \langle \alpha, \beta' \rangle\)，对于固定的 \(\alpha\)，是 \(\alpha\) 相对于基 \((\overline{\omega}_{\beta})_{\beta \in B}\) 的坐标：

\[
\alpha = \sum_ {\beta \in B} n (\alpha , \beta) \overline {{\omega}} _ {\beta}.\tag{14}
\]

因此，Cartan 矩阵是典范嵌入的矩阵的转置

\[
\mathrm{Q} (\mathrm{R}) \rightarrow \mathrm{P} (\mathrm{R})
\]

相对于基 B 与 \((\overline{\omega}_{\alpha})_{\alpha\in\mathsf{B}}\) 的 Z-模 \(\mathrm{Q}(\mathrm{R})\) 与 \(\mathrm{P}(\mathrm{R})\)，权 \(\overline{\omega}\) 若属于 \(\overline{C}\)，便称为支配权；换言之，它相对于 \((\overline{\omega}_{\alpha})_{\alpha\in\mathsf{B}}\) 的坐标都是 \(\geqslant0\) 的整数，或者等价地，若 \(g(\overline{\omega})\leqslant\overline{\omega}\) 对所有 \(g\in\mathrm{W}(\mathrm{R})\) 成立。由于 \(\overline{C}\) 是 \(\mathrm{W(R)}\) 的基本域，对于任意权 \(\overline{\omega}\)，存在唯一的权 \(\overline{\omega}'\)，使得 \(\overline{\omega}'\) 是 \(\overline{\omega}\) 在 \(\mathrm{W(R)}\) 下的变换。

有

\[
\langle \overline {{{{\omega}}}} _ {\alpha}, \beta^ {\prime} \rangle = (\overline {{{{\omega}}}} _ {\alpha} | \frac {2 \beta}{(\beta | \beta)}) = \delta_ {\alpha \beta}
\]

对于 \(\alpha, \beta \in \mathbf{B}\)（其中 \(\delta_{\alpha\beta}\) 表示 Kronecker 符号），于是

\[
s _ {\beta} (\overline {{\omega}} _ {\alpha}) = \overline {{\omega}} _ {\alpha} - \delta_ {\alpha \beta} \beta \quad \text { and } \quad (\overline {{\omega}} _ {\alpha} | \beta) = \frac {1}{2} (\beta | \beta) \delta_ {\alpha \beta}.
\]

 换言之，\(\overline{\omega}_{\alpha}\) 与 \(\beta\) 正交（当 \(\beta \neq \alpha\) 时），而它在 \(R\alpha\) 上的正交投影是 \(\frac{1}{2}\alpha\)。由于 \(\overline{\omega}_{\alpha} \in \overline{C}\)，有 \((\overline{\omega}_{\alpha} | \overline{\omega}_{\beta}) \geqslant 0\) 对 \(\alpha, \beta \in B\) 成立，即角 \((\widehat{\overline{\omega}_{\alpha}, \overline{\omega}_{\beta}})\) 是锐角或直角。支配权是满足下述条件的 \(\overline{\omega} \in V\)：其 \(2(\overline{\omega} | \alpha)/( \alpha | \alpha)\) 是一个整数 \(\geqslant 0\)，对所有 \(\alpha \in B\) 都成立。

命题 28. 设 B 是 R 的一个基，B' 是 B 的子集，V' 是由 B' 生成的 V 的向量子空间，且 R' = R ∩ V'（它是 V' 中的根系），R' 是逆根系（它与 R' 在 R' 中的典范像等同），V1 是 R' 在 V 中的正交补，p 是沿 V1 将 V 投影到 V' 的投影。于是，Q(R') = Q(R) ∩ V'，P(R') = p(P(R))。R' 的支配权集合是 R 的支配权集合在 p 下的像。

事实上，\(\mathrm{Q}(\mathrm{R})\) 是 \(\mathrm{V}\) 的子群，以 \(\mathrm{B}\) 为基；\(\mathrm{Q}(\mathrm{R}')\) 是 \(\mathrm{V}'\) 的子群，以 \(\mathrm{B}'\) 为基（no. 7, Cor. 4 of Prop. 20），由此 \(\mathrm{Q}(\mathrm{R}') = \mathrm{Q}(\mathrm{R}) \cap \mathrm{V}'\) 立即可得。若 \(\overline{\omega} \in \mathrm{P}(\mathrm{R})\) 且 \(\alpha \in \mathrm{R}'\)，则 \(\langle p(\overline{\omega}), \alpha' \rangle = \langle \overline{\omega}, \alpha' \rangle \in \mathbf{Z}\)，所以 \(p(\overline{\omega}) \in \mathrm{P}(\mathrm{R}')\)，从而 \(p(\mathrm{P}(\mathrm{R})) \subseteq \mathrm{P}(\mathrm{R}')\)。若 \(\overline{\omega}' \in \mathrm{P}(\mathrm{R}')\)，则 \(\overline{\omega}'\) 延拓为线性形式 \(\overline{\omega}\)，定义在 \(\mathrm{V}^*\) 上，并在 \((\mathrm{B} - \mathrm{B}')'\) 上为零；有 \(\langle \overline{\omega}, \alpha' \rangle \in \mathbf{Z}\) 对所有 \(\alpha \in \mathrm{B}\) 成立，所以 \(\overline{\omega} \in \mathrm{P}(\mathrm{R})\)，且 \(\overline{\omega}' = p(\overline{\omega})\)；于是 \(\mathrm{P}(\mathrm{R}') \subseteq p(\mathrm{P}(\mathrm{R}))\)。因此 \(\mathrm{P}(\mathrm{R}') = p(\mathrm{P}(\mathrm{R}))\)，而关于支配权的断言以同样方式得到。

命题 29. 设 \(\rho\) 是所有 \(>0\) 的根之和的一半。

\begin{enumerate}
\def\labelenumi{(\roman{enumi})}
\item
  \(\rho = \sum_{\alpha \in B} \overline{\omega}_{\alpha}\)；这是 C 中的元素。
\item
  有 \(s_{\alpha}(\rho) = \rho -\alpha\) 对所有 \(\alpha \in B\) 成立
\item
  \((2\rho |\alpha) = (\alpha |\alpha)\) 对所有 \(\alpha \in B\) 成立。
\end{enumerate}

由于 R 是约化的，\(s_{\alpha}(R_{+}(C) - \{\alpha\}) = R_{+}(C) - \{\alpha\}\) 且 \(s_{\alpha}(\alpha) = -\alpha\) 对所有 \(\alpha \in B\) 成立（no. 6, Cor. 1 of Prop. 17），所以 \(s_{\alpha}(2\rho) = 2\rho - 2\alpha\)。由于 \(s_{\alpha}(\rho) = \rho - \langle \rho, \alpha^{-} \rangle\)，我们看到

\[
\langle \rho , \alpha^ {\prime} \rangle = 1 = \left\langle \sum_ {\beta \in B} \bar {\omega} _ {\beta}, \alpha^ {\prime} \right\rangle .
\]

于是，\(\rho = \sum_{\beta} \overline{\omega}_{\beta}\)，从而 \(\rho \in C\)。最后，(iii) 等价于 \(\langle \rho, \alpha^{-} \rangle = 1\)。

推论。设 \(\sigma\) 是所有 \(>0\) 的 \(\mathbf{R}^{-}\) 元素之和的一半（相对于 \(\mathbf{B}^{-}\)）。对所有 \(\alpha \in V\)，\(\alpha\) 关于基 \(\mathbf{B}\) 的坐标之和是 \(\langle \alpha, \sigma \rangle\)。若 \(\alpha \in R\)，这个和等于 \(\frac{1}{2} \sum_{\beta \in \mathbb{R}_{+}(C)} n(\alpha, \beta)\)。

交换上文 R 与 \(R^{-}\) 的角色，有 \(\langle\alpha,\sigma\rangle=1\) 对所有 \(\alpha\in B\) 成立，故推论得证。

